% GENERATED FILE. Edit canonical lessons, then run npm run generate:books.
% canonical-source-hash: fnv1a64:40645ec73e848662
% canonical-lessons: UR-C503-silan, UR-C503-int, UR-C503-palastar, UR-C503-rangvana, UR-C503-lipna

\chapter{Damp and Cracks}
\label{ch:ur-damp-and-cracks}
\hlchaptermodality{503}

\begin{chapteropening}
\textbf{By the end of this chapter:} \emph{I can talk about damp walls and painting them.}

Complete the last lesson of chapter 503: I can talk about damp walls and painting them.
\end{chapteropening}

\section[\texorpdfstring{\ur{سیلن} (sīlan)}{سیلن (sīlan)}]{\ur{سیلن} (sīlan) — damp (in a wall)}

\label{lesson:UR-C503-silan}



\begin{quote}
Before the new one: say the Urdu for rust, then the Urdu for a metal spring.
\end{quote}

\subsection*{You'll want to know: \ur{سیلن}}
\textbf{\ur{سیلن}} — \emph{sīlan} — \textquotedblleft{}damp (in a wall)\textquotedblright{}.

\begin{quote}
\ur{برسات} \ur{میں} \ur{دیواروں} \ur{میں} \ur{سیلن} \ur{آ} \ur{جاتی} \ur{ہے۔}
\end{quote}

\begin{quote}
\emph{barsāt meṅ dīvāroṅ meṅ sīlan ā jātī hai.}
\end{quote}

\begin{quote}
In the rainy season damp gets into the walls.
\end{quote}

Feminine.

\begin{cousinweb}
\textbf{\ur{سیلن}} (\emph{sīlan}) is built on \emph{sīl}, \textquotedblleft{}cool, moist\textquotedblright{}, which CDIAL lists under Sanskrit \emph{sītala}, \textquotedblleft{}cold\textquotedblright{} (CDIAL 12487).
\end{cousinweb}

\begin{grammarlens}[title={What you've built}]
One more word for describing what you lived and felt.
\end{grammarlens}

\subsection*{Your turn}
\begin{itemize}
\raggedright
  \item \emph{Say it:} \emph{sīlan}
  \item \emph{Say it:} \emph{sīlan}, once more
  \item \emph{Say it:} \emph{barsāt meṅ dīvāroṅ meṅ sīlan ā jātī hai.} — In the rainy season damp gets into the walls.
  \item \emph{Recall:} say the Urdu for rust, then the Urdu for a metal spring, then say \emph{sīlan} again
\end{itemize}

\begin{tcolorbox}[breakable,colback=teal!4,colframe=teal!35!black,title={Before you move on}]
What does \ur{سیلن} mean? (\textquotedblleft{}Damp (in a wall)\textquotedblright{}.) Say it once more.
\end{tcolorbox}
\section[\texorpdfstring{\ur{اینٹ} (īṅṭ)}{اینٹ (īṅṭ)}]{\ur{اینٹ} (īṅṭ) — a brick}

\label{lesson:UR-C503-int}



\begin{quote}
Before the new one: say the Urdu for a metal spring, then the Urdu for damp (in a wall).
\end{quote}

\subsection*{You'll want to know: \ur{اینٹ}}
\textbf{\ur{اینٹ}} — \emph{īṅṭ} — \textquotedblleft{}a brick\textquotedblright{}.

\begin{quote}
\ur{مزدور} \ur{نے} \ur{دیوار} \ur{کے} \ur{لیے} \ur{اینٹیں} \ur{اٹھائیں۔}
\end{quote}

\begin{quote}
\emph{mazdūr ne dīvār ke liye īṅṭeṅ uṭhāīṅ.}
\end{quote}

\begin{quote}
The labourer carried bricks for the wall.
\end{quote}

Feminine.

\begin{cousinweb}
\textbf{\ur{اینٹ}} (\emph{īṅṭ}) is inherited from Sanskrit \emph{iṣṭakā}, \textquotedblleft{}a brick\textquotedblright{} (CDIAL 1600).
\end{cousinweb}

\begin{grammarlens}[title={What you've built}]
One more word for describing what you lived and felt.
\end{grammarlens}

\subsection*{Your turn}
\begin{itemize}
\raggedright
  \item \emph{Say it:} \emph{īṅṭ}
  \item \emph{Say it:} \emph{īṅṭ}, once more
  \item \emph{Say it:} \emph{mazdūr ne dīvār ke liye īṅṭeṅ uṭhāīṅ.} — The labourer carried bricks for the wall.
  \item \emph{Recall:} say the Urdu for a metal spring, then the Urdu for damp (in a wall), then say \emph{īṅṭ} again
\end{itemize}

\begin{tcolorbox}[breakable,colback=teal!4,colframe=teal!35!black,title={Before you move on}]
What does \ur{اینٹ} mean? (\textquotedblleft{}A brick\textquotedblright{}.) Say it once more.
\end{tcolorbox}
\section[\texorpdfstring{\ur{پلستر} (palastar)}{پلستر (palastar)}]{\ur{پلستر} (palastar) — plaster (on a wall)}

\label{lesson:UR-C503-palastar}



\begin{quote}
Before the new one: say the Urdu for damp (in a wall), then the Urdu for a brick.
\end{quote}

\subsection*{You'll want to know: \ur{پلستر}}
\textbf{\ur{پلستر}} — \emph{palastar} — \textquotedblleft{}plaster (on a wall)\textquotedblright{}.

\begin{quote}
\ur{برسات} \ur{میں} \ur{دیوار} \ur{کا} \ur{پلستر} \ur{گر} \ur{گیا۔}
\end{quote}

\begin{quote}
\emph{barsāt meṅ dīvār kā palastar gir gayā.}
\end{quote}

\begin{quote}
The plaster came off the wall in the rainy season.
\end{quote}

Masculine.

\begin{cousinweb}
\textbf{\ur{پلستر}} (\emph{palastar}) is borrowed from English \emph{plaster}.
\end{cousinweb}

\begin{grammarlens}[title={What you've built}]
One more word for describing what you lived and felt.
\end{grammarlens}

\subsection*{Your turn}
\begin{itemize}
\raggedright
  \item \emph{Say it:} \emph{palastar}
  \item \emph{Say it:} \emph{palastar}, once more
  \item \emph{Say it:} \emph{barsāt meṅ dīvār kā palastar gir gayā.} — The plaster came off the wall in the rainy season.
  \item \emph{Recall:} say the Urdu for damp (in a wall), then the Urdu for a brick, then say \emph{palastar} again
\end{itemize}

\begin{tcolorbox}[breakable,colback=teal!4,colframe=teal!35!black,title={Before you move on}]
What does \ur{پلستر} mean? (\textquotedblleft{}Plaster (on a wall)\textquotedblright{}.) Say it once more.
\end{tcolorbox}
\section[\texorpdfstring{\ur{رنگوانا} (rangvānā)}{رنگوانا (rangvānā)}]{\ur{رنگوانا} (rangvānā) — to have (something) painted}

\label{lesson:UR-C503-rangvana}



\begin{quote}
Before the new one: say the Urdu for a brick, then the Urdu for plaster (on a wall).
\end{quote}

\subsection*{You'll want to know: \ur{رنگوانا}}
\textbf{\ur{رنگوانا}} — \emph{rangvānā} — \textquotedblleft{}to have (something) painted\textquotedblright{}.

\begin{quote}
\ur{ہم} \ur{نے} \ur{برسات} \ur{کے} \ur{بعد} \ur{دیواریں} \ur{رنگوائیں۔}
\end{quote}

\begin{quote}
\emph{ham ne barsāt ke bād dīvāreṅ rangvāīṅ.}
\end{quote}

\begin{quote}
We had the walls painted after the rainy season.
\end{quote}

Past, with \ur{نے} (he or she …): \ur{اس} \ur{نے} \ur{دروازہ} \ur{رنگوایا}, \emph{us ne darvāzā rangvāyā}.

\begin{cousinweb}
\textbf{\ur{رنگوانا}} (\emph{rangvānā}) is the causative of \emph{rangnā}, \textquotedblleft{}to colour\textquotedblright{}: to have someone colour a thing. CDIAL lists the Hindi cousin \emph{rãgānā}, \textquotedblleft{}to be coloured\textquotedblright{}, under \emph{*raṅgayati}, \textquotedblleft{}dyes\textquotedblright{}, a form it reconstructs (CDIAL 10570).
\end{cousinweb}

\begin{grammarlens}[title={What you've built}]
One more word for describing what you lived and felt.
\end{grammarlens}

\subsection*{Your turn}
\begin{itemize}
\raggedright
  \item \emph{Say it:} \emph{rangvānā}
  \item \emph{Say it:} \emph{rangvānā}, once more
  \item \emph{Say it:} \emph{ham ne barsāt ke bād dīvāreṅ rangvāīṅ.} — We had the walls painted after the rainy season.
  \item \emph{Recall:} say the Urdu for a brick, then the Urdu for plaster (on a wall), then say \emph{rangvānā} again
\end{itemize}

\begin{tcolorbox}[breakable,colback=teal!4,colframe=teal!35!black,title={Before you move on}]
What does \ur{رنگوانا} mean? (\textquotedblleft{}To have (something) painted\textquotedblright{}.) Say it once more.
\end{tcolorbox}
\section[\texorpdfstring{\ur{لیپنا} (līpnā)}{لیپنا (līpnā)}]{\ur{لیپنا} (līpnā) — to smear (with mud)}

\label{lesson:UR-C503-lipna}



\begin{quote}
Before the new one: say the Urdu for plaster (on a wall), then the Urdu for to have (something) painted.
\end{quote}

\subsection*{You'll want to know: \ur{لیپنا}}
\textbf{\ur{لیپنا}} — \emph{līpnā} — \textquotedblleft{}to smear (with mud)\textquotedblright{}.

\begin{quote}
\ur{دادی} \ur{نے} \ur{آنگن} \ur{کو} \ur{مٹی} \ur{سے} \ur{لیپا۔}
\end{quote}

\begin{quote}
\emph{dādī ne āṅgan ko miṭṭī se līpā.}
\end{quote}

\begin{quote}
Grandmother smeared the courtyard with fresh mud.
\end{quote}

Past, with \ur{نے} (he or she …): \ur{اس} \ur{نے} \ur{فرش} \ur{لیپا}, \emph{us ne farsh līpā}.

\begin{cousinweb}
\textbf{\ur{لیپنا}} (\emph{līpnā}) is inherited: CDIAL lists the Hindi cousin \emph{līpnā}, \textquotedblleft{}to smear\textquotedblright{}, under Sanskrit \emph{lipyatē}, \textquotedblleft{}is smeared\textquotedblright{} (CDIAL 11061).
\end{cousinweb}

\begin{grammarlens}[title={What you've built}]
One more word for describing what you lived and felt.
\end{grammarlens}

\subsection*{Your turn}
\begin{itemize}
\raggedright
  \item \emph{Say it:} \emph{līpnā}
  \item \emph{Say it:} \emph{līpnā}, once more
  \item \emph{Say it:} \emph{dādī ne āṅgan ko miṭṭī se līpā.} — Grandmother smeared the courtyard with fresh mud.
  \item \emph{Recall:} say the Urdu for plaster (on a wall), then the Urdu for to have (something) painted, then say \emph{līpnā} again
\end{itemize}

\begin{tcolorbox}[breakable,colback=teal!4,colframe=teal!35!black,title={Before you move on}]
What does \ur{لیپنا} mean? (\textquotedblleft{}To smear (with mud)\textquotedblright{}.) Say it once more.
\end{tcolorbox}
